% 分类目录：ml
\documentclass[12pt, a4paper]{article}
\usepackage[margin=2.5cm]{geometry}
\usepackage{ctex}
\usepackage{amsmath, amssymb}
\usepackage{listings}
\usepackage{xcolor}
\usepackage{tabularx}
\usepackage{hyperref}

\lstset{
    language=Python,
    basicstyle=\ttfamily\small,
    keywordstyle=\color{blue},
    commentstyle=\color{green!50!black},
    stringstyle=\color{red},
    showstringspaces=false,
    breaklines=true,
    numbers=left,
    numberstyle=\tiny\color{gray},
    frame=single
}

\title{知识点：概率论基础 (Probability Theory Foundations)}
\author{}
\date{}

\begin{document}
\maketitle

\section{核心对象}
概率论是研究随机现象数量规律的数学分支。在机器学习中，它不仅用于建模数据的不确定性，还被用来推导损失函数（如极大似然估计）。
\begin{itemize}
    \item \textbf{随机变量 (Random Variable)}：其值随随机实验结果而定的变量，分为离散型和连续型。
    \item \textbf{概率质量函数 (PMF)}：离散变量在各取值点的概率。
    \item \textbf{概率密度函数 (PDF)}：连续型变量在其取值范围内各点附近的概率分布率。
    \item \textbf{累积分布函数 (CDF)}：变量小于或等于某值的总概率。
\end{itemize}

\section{离散型概率分布 (Discrete Distributions)}

\subsection{伯努利分布 (Bernoulli)}
\textbf{描述}：单次试验，只有两种结果（成功或失败）。
\\ \textbf{PMF}：$P(X=k) = p^k(1-p)^{1-k}, \quad k \in \{0, 1\}$
\\ \textbf{特征}：$E[X] = p$， $Var(X) = p(1-p)$
\\ \textbf{ML 应用}：逻辑回归的输出建模。

\subsection{二项分布 (Binomial)}
\textbf{描述}：$n$ 次独立的伯努利试验中成功的次数。
\\ \textbf{PMF}：$P(X=k) = \binom{n}{k} p^k(1-p)^{n-k}, \quad k \in \{0, 1, \dots, n\}$
\\ \textbf{特征}：$E[X] = np$， $Var(X) = np(1-p)$

\subsection{几何分布 (Geometric)}
\textbf{描述}：在伯努利试验中，首次成功出现在第 $k$ 次试验的概率。
\\ \textbf{PMF}：$P(X=k) = (1-p)^{k-1}p, \quad k \in \{1, 2, \dots\}$
\\ \textbf{特征}：$E[X] = \frac{1}{p}$， $Var(X) = \frac{1-p}{p^2}$

\subsection{泊松分布 (Poisson)}
\textbf{描述}：在单位时间或空间内，某稀有事件发生的次数。
\\ \textbf{PMF}：$P(X=k) = \frac{\lambda^k e^{-\lambda}}{k!}, \quad k \in \{0, 1, \dots\}$
\\ \textbf{特征}：$E[X] = \lambda$， $Var(X) = \lambda$

\subsection{负二项分布 (Negative Binomial)}
\textbf{描述}：在伯努利试验中，第 $r$ 次成功发生时恰好经历了 $k$ 次失败。
\\ \textbf{PMF}：$P(X=k) = \binom{k+r-1}{k} p^r(1-p)^k$
\\ \textbf{特征}：$E[X] = \frac{r(1-p)}{p}$

\subsection{超几何分布 (Hypergeometric)}
\textbf{描述}：不放回抽样模型。在含有 $M$ 个次品、共 $N$ 个物品中抽取 $n$ 个，抽到 $k$ 个次品的概率。
\\ \textbf{PMF}：$P(X=k) = \frac{\binom{M}{k}\binom{N-M}{n-k}}{\binom{N}{n}}$

\subsection{多项分布 (Multinomial)}
\textbf{描述}：二项分布的推广。$n$ 次试验，每次试验有 $k$ 种可能结果。
\\ \textbf{PMF}：$P(X_1=x_1, \dots, X_k=x_k) = \frac{n!}{x_1! \dots x_k!} p_1^{x_1} \dots p_k^{x_k}$
\\ \textbf{ML 应用}：分类任务中的分类标签建模。

\section{连续型一元分布 (Continuous Univariate Distributions)}

\subsection{均匀分布 (Uniform)}
\textbf{PDF}：$f(x) = \frac{1}{b-a}, \quad a \le x \le b$
\\ \textbf{特征}：$E[X] = \frac{a+b}{2}, \quad Var(X) = \frac{(b-a)^2}{12}$

\subsection{正态分布 (Normal / Gaussian)}
\textbf{PDF}：$f(x) = \frac{1}{\sigma\sqrt{2\pi}} e^{-\frac{(x-\mu)^2}{2\sigma^2}}$
\\ \textbf{特征}：$E[X] = \mu, \quad Var(X) = \sigma^2$
\\ \textbf{ML 应用}：最小二乘法的噪声假设，深度学习中的权重初始化。

\subsection{指数分布 (Exponential)}
\textbf{PDF}：$f(x) = \lambda e^{-\lambda x}, \quad x \ge 0$
\\ \textbf{CDF}：$F(x) = 1 - e^{-\lambda x}$
\\ \textbf{特征}：$E[X] = \frac{1}{\lambda}, \quad Var(X) = \frac{1}{\lambda^2}$

\subsection{伽马分布 (Gamma)}
\textbf{PDF}：$f(x) = \frac{\beta^\alpha}{\Gamma(\alpha)} x^{\alpha-1} e^{-\beta x}, \quad x > 0$
\\ \textbf{特征}：$E[X] = \frac{\alpha}{\beta}, \quad Var(X) = \frac{\alpha}{\beta^2}$

\subsection{贝塔分布 (Beta)}
\textbf{描述}：定义在 $[0, 1]$ 上的分布，常作伯努利参数 $p$ 的先验。
\\ \textbf{PDF}：$f(x) = \frac{x^{\alpha-1}(1-x)^{\beta-1}}{B(\alpha, \beta)}$
\\ \textbf{特征}：$E[X] = \frac{\alpha}{\alpha+\beta}$

\subsection{卡方分布 (Chi-square)}
\textbf{描述}：$k$ 个独立标准正态随机变量的平方和。
\\ \textbf{特征}：$E[X] = k, \quad Var(X) = 2k$

\subsection{t 分布 (Student's t)}
\textbf{描述}：小样本统计推断，具有比正态分布更厚的尾部（Heavy Tails）。
\\ \textbf{ML 应用}：异常检测。

\subsection{F 分布 (F-distribution)}
\textbf{描述}：两个独立卡方分布除以各自自由度后的比值。
\\ \textbf{应用}：方差分析 (ANOVA)。

\subsection{对数正态分布 (Log-normal)}
\textbf{描述}：若 $\ln X \sim \mathcal{N}(\mu, \sigma^2)$，则 $X$ 服从对数正态分布。
\\ \textbf{PDF}：$f(x) = \frac{1}{x\sigma\sqrt{2\pi}} e^{-\frac{(\ln x-\mu)^2}{2\sigma^2}}$

\subsection{帕累托分布 (Pareto)}
\textbf{描述}：幂律分布，“二八原则”的体现。
\\ \textbf{PDF}：$f(x) = \frac{\alpha x_m^\alpha}{x^{\alpha+1}}, \quad x \ge x_m$

\section{多元概率分布 (Multivariate Distributions)}

\subsection{多元正态分布 (Multivariate Normal)}
\textbf{PDF}：$f(\mathbf{x}) = \frac{1}{(2\pi)^{k/2}|\boldsymbol{\Sigma}|^{1/2}} \exp\left(-\frac{1}{2}(\mathbf{x}-\boldsymbol{\mu})^T \boldsymbol{\Sigma}^{-1} (\mathbf{x}-\boldsymbol{\mu})\right)$
\\ \textbf{关键参数}：均值向量 $\boldsymbol{\mu}$，协方差矩阵 $\boldsymbol{\Sigma}$。

\subsection{狄利克雷分布 (Dirichlet)}
\textbf{描述}：多项分布对应的共轭先验分布，取值在单纯形（Simplex）上。
\\ \textbf{ML 应用}：主题模型 (LDA)。

\subsection{威沙特分布 (Wishart)}
\textbf{描述}：多元正态分布中“样本协方差矩阵”的分布，是多变量统计分析的重要基础。

\section{信息论基础}
\begin{itemize}
    \item \textbf{熵 (Entropy)}：衡量信息的不确定性，$H(X) = - \sum p(x) \log p(x)$。
    \item \textbf{交叉熵 (Cross Entropy)}：评估两分布的一致性，分类任务中最常用的损失函数。
\end{itemize}

\section{深度面试题}

\textbf{问：为什么贝塔分布是伯努利分布的共轭先验？}\\
答：共轭先验意味着先验分布与后验分布属于同一种分布。伯努利似然的形式是 $p^k(1-p)^{n-k}$，而贝塔先验的形式是 $p^{\alpha-1}(1-p)^{\beta-1}$。相乘后，指数项相加，所得后验依然是一个贝塔分布，只是参数变为了 $(\alpha+k, \beta+n-k)$。这使得参数更新非常直观。

\section{参考文献}
\begin{enumerate}
    \item Ross, S. M. (2014). \textit{A First Course in Probability}. Pearson.
    \item Goodfellow, I., Bengio, Y., \& Courville, A. (2016). \textit{Deep Learning}. MIT Press (Chapter 3).
\end{enumerate}

\end{document}
